\documentclass[notes,10pt, aspectratio=169]{beamer}

\usepackage{pgfpages}
% These slides also contain speaker notes. You can print just the slides,
% just the notes, or both, depending on the setting below. Comment out the want
% you want.
% Note visibility is selected after package loading at the end of this preamble.
%\setbeameroption{show only notes} % Only notes
%\setbeameroption{show notes on second screen=right} % Both

\usepackage{helvet}
\usepackage[default]{lato}
\usepackage{array}
\usepackage{eurosym}
\usepackage{setspace}
\usepackage{animate}
\usepackage{multicol}
\usepackage{breqn}
\usepackage{subcaption}
\usepackage{caption}
\captionsetup{font=scriptsize, justification=raggedright,singlelinecheck=false}
\usepackage[export]{adjustbox}
\usepackage[scale=2]{ccicons}
\usepackage{tikz}
\newenvironment{wideitemize}{\itemize\addtolength{\itemsep}{10pt}}{\enditemize}
\usepackage{natbib}

\usepackage{tikz}
\usepackage{verbatim}
\setbeamertemplate{note page}{\pagecolor{yellow!5}\insertnote}
\usepackage{amsmath}
\usepackage{bm}
%\usepackage{mathpazo}
\usepackage{hyperref}
\usepackage{lipsum}
\usepackage{multimedia}
\usepackage{multirow}
\usepackage{graphicx}
\usepackage{dcolumn}
\usepackage{bbm}
\newcolumntype{d}[0]{D{.}{.}{5}}

\usepackage{changepage}
\usepackage{appendixnumberbeamer}
% \newcommand{\beginbackup}{
%   \newcounter{framenumbervorappendix}
%   \setcounter{framenumbervorappendix}{\value{framenumber}}
%   \setbeamertemplate{footline}
%   {
%      \leavevmode%
%      \hline
%      box{%
%       \begin{beamercolorbox}[wd=\paperwidth,ht=2.25ex,dp=1ex,right]{footlinecolor}%
% %         \insertframenumber  \hspace*{2ex}
%       \end{beamercolorbox}}%
%      \vskip0pt%
%   }
%  }

 \setbeamertemplate{footline}{
    \hbox{%
    \begin{beamercolorbox}[wd=\paperwidth,ht=3ex,dp=1.5ex,leftskip=2ex,rightskip=2ex]{page footer}%
        %\usebeamerfont{title in head/foot}%
        %\insertshorttitle \hfill
        %\insertsection \hfill
        %\tableofcontents[currentsection] \hfill
        \insertsectionnavigationhorizontal{0.7\paperwidth}{}{} \hfill
        \insertframenumber{} / \inserttotalframenumber
    \end{beamercolorbox}}%
}
\newcommand{\backupend}{
   \addtocounter{framenumbervorappendix}{-\value{framenumber}}
   \addtocounter{framenumber}{\value{framenumbervorappendix}}
}


\usepackage{graphicx}
\usepackage[space]{grffile}
\usepackage{booktabs}

% These are my colors -- there are many like them, but these ones are mine.
\definecolor{blue}{RGB}{0,114,178}
\definecolor{red}{RGB}{213,94,0}
\definecolor{yellow}{RGB}{240,228,66}
\definecolor{green}{RGB}{0,158,115}

\setbeamercolor{section in head/foot}{fg=blue}

\useinnertheme[shadow]{rounded}
\setbeamercolor{block title example}{fg=black,bg=orange!50!white}
\setbeamercolor{block body example}{fg=black, bg=orange!30!white}


\hypersetup{
  colorlinks=false,
  linkbordercolor = {white},
  linkcolor = {blue}
}


%% I use a beige off white for my background
\definecolor{MyBackground}{RGB}{255,253,218}

%% Uncomment this if you want to change the background color to something else
%\setbeamercolor{background canvas}{bg=MyBackground}

%% Change the bg color to adjust your transition slide background color!
\newenvironment{transitionframe}{
  \setbeamercolor{background canvas}{bg=yellow}
  \begin{frame}}{
    \end{frame}
}

\setbeamercolor{frametitle}{fg=blue}
\setbeamercolor{title}{fg=black}
%\setbeamertemplate{footline}[frame number]
\setbeamertemplate{navigation symbols}{}
\setbeamertemplate{itemize items}{-}
\setbeamercolor{itemize item}{fg=blue}
\setbeamercolor{itemize subitem}{fg=blue}
\setbeamercolor{enumerate item}{fg=blue}
\setbeamercolor{enumerate subitem}{fg=blue}
\setbeamercolor{button}{bg=white,fg=blue,}



% If you like road maps, rather than having clutter at the top, have a roadmap show up at the end of each section
% (and after your introduction)
% Uncomment this is if you want the roadmap!
% \AtBeginSection[]
% {
%    \begin{frame}
%        \frametitle{Roadmap of Talk}
%        \tableofcontents[currentsection]
%    \end{frame}
% }
\setbeamercolor{section in toc}{fg=blue}
\setbeamercolor{subsection in toc}{fg=red}
%\setbeamersize{text margin left=1em,text margin right=1em}
\setbeamersize{text margin left=0.7cm,text margin right=0.7cm}


\title[]{\textcolor{blue}{Seminar 6:\\Inference and Interpretation}}
\subtitle{Five problems and worked solutions}
\author{25117 - Econometrics}
\institute[shortinst]{Universitat Pompeu Fabra}

%\date{September 25th, 2024}


% Each problem is followed by worked solution slides. Set \showsolutionsfalse
% below to compile a problem-only deck. All data are hypothetical;
% no external assets are needed. Supplied uncertainty estimates support
% large-sample inference; no covariance-estimation formula is required.
\newif\ifshowsolutions
\showsolutionstrue
\newcommand{\E}{\mathbb{E}}
\newcommand{\Var}{\operatorname{Var}}
\newcommand{\Cov}{\operatorname{Cov}}
\newcommand{\Corr}{\operatorname{Corr}}
\newcommand{\answerpart}[1]{\textcolor{blue}{\textbf{#1.}}\,}

\setbeameroption{hide notes} % Speaker notes hidden; worked solutions are visible.


\date{}
\begin{document}
\begin{frame}[plain]
\maketitle
\end{frame}

\section{Inference}
\begin{frame}{Problem 1: Is class size associated with test scores?}
\small
A hypothetical sample of 500 schools gives
\[
\widehat{Score}_i=620-1.50\,ClassSize_i+2.00\,Income_i,
\]
where scores are in points, class size is students per class, and income is in thousands of euros. The supplied standard error of the class-size coefficient is $0.60$.
\par\smallskip
Use the large-sample normal approximation, a 5\% significance level, and $\Phi(2.50)=0.9938$.
\begin{enumerate}
\item Test $H_0:\beta_{ClassSize}=0$ against a two-sided alternative. Calculate the statistic, approximate $p$-value, and 95\% confidence interval.
\item Test $H_0:\beta_{ClassSize}=-2$ using the same estimate and standard error. Check the decision against your interval.
\item Interpret the estimated coefficient, holding income fixed. Does statistical significance establish that reducing class size causes higher scores?
\end{enumerate}
\end{frame}

\ifshowsolutions
\begin{frame}{Solution 1: Tests, confidence interval, and interpretation}
\small
\answerpart{1}For the zero null,
\[
t=\frac{-1.50}{0.60}=-2.50,\text{ and }p\approx2[1-\Phi(2.50)]=0.0124.
\]
Reject at 5\%, since $|t|>1.96$ and $p<0.05$. The interval is
\[
-1.50\pm1.96(0.60)=[-2.676,-0.324].
\]
\answerpart{2}For the null value $-2$,
\[
t=\frac{-1.50-(-2)}{0.60}=0.8333.
\]
Do not reject: $|t|<1.96$, and $-2$ lies inside the 95\% interval.
\par\smallskip
\answerpart{3}An additional student per class is associated with a 1.50-point lower predicted score, holding district income fixed. The test concerns the regression coefficient. A causal interpretation also requires credible control of confounding; significance alone does not establish causality.
\end{frame}
\fi

\begin{frame}{Problem 2: Can two insignificant coefficients matter jointly?}
\small
A hypothetical regression of test scores on weekly tutoring hours $X_1$, weekly independent-study hours $X_2$, and background controls uses $n=600$ students. For the two slopes, software supplies
\[
\hat\beta_1=1.50,\text{ and }\hat\beta_2=1.50,\text{ with }
\widehat V=\begin{pmatrix}1&-0.50\\-0.50&1\end{pmatrix}.
\]
Rows and columns of $\widehat V$ follow $(\hat\beta_1,\hat\beta_2)$. Use the supplied covariance estimate and large-sample inference.
\begin{enumerate}
\item Read both standard errors and the estimated correlation $\hat\rho$. Test each coefficient separately against zero at 5\%.
\item State the joint null and alternative, count the restrictions, and calculate
\[
F=\frac{t_1^2+t_2^2-2\hat\rho t_1t_2}{2(1-\hat\rho^2)}.
\]
\item The large-sample 5\% critical value for $q=2$ is $3.00$. Give the joint decision. Explain why averaging the two squared $t$ statistics is incorrect here.
\end{enumerate}
\end{frame}

\begin{frame}{The joint test: General variances and matrix form}
\footnotesize
\setlength{\abovedisplayskip}{3pt}
\setlength{\belowdisplayskip}{3pt}
Test $H_0:\beta_1=\beta_{1,0},\ \beta_2=\beta_{2,0}$ ($q=2$ restrictions). Define
\[
d_j=\hat\beta_j-\beta_{j,0},\text{ }s_j=SE(\hat\beta_j),
\text{ and }c=\widehat{\Cov}(\hat\beta_1,\hat\beta_2).
\]
\textbf{Matrix form.} Collect the deviations and their estimated covariance matrix:
\[
\boldsymbol d=\begin{pmatrix}d_1\\d_2\end{pmatrix},
\text{ }\widehat V=\begin{pmatrix}s_1^2&c\\c&s_2^2\end{pmatrix},
\text{ and }\boxed{F=\frac{1}{2}\boldsymbol d'\widehat V^{-1}\boldsymbol d}.
\]
Here $'$ denotes transpose and $^{-1}$ denotes matrix inverse. Assume $s_1^2s_2^2>c^2$.
\par\smallskip
\textbf{Expanded form.} Inverting this $2\times2$ matrix gives
\[
\boxed{F=\frac{s_2^2d_1^2+s_1^2d_2^2-2c\,d_1d_2}{2(s_1^2s_2^2-c^2)}}.
\]
\textbf{Standardized form.} Substitute $t_j=d_j/s_j$ and $\hat\rho=c/(s_1s_2)$:
\[
F=\frac{t_1^2+t_2^2-2\hat\rho t_1t_2}{2(1-\hat\rho^2)}.
\]
All three forms allow unequal variances. In Problem 2, $s_1=s_2=1$ simply makes $t_j=d_j$.
\end{frame}

\ifshowsolutions
\begin{frame}{Solution 2: The covariance changes the joint test}
\small
\answerpart{1}Standard errors are square roots of the diagonal entries:
\[
SE(\hat\beta_1)=SE(\hat\beta_2)=1,\text{ and }\hat\rho=\frac{-0.50}{1\cdot1}=-0.50.
\]
Thus $t_1=t_2=1.50$. Neither separate test rejects at 5\%.
\par\smallskip
\answerpart{2}The joint null is $H_0:\beta_1=\beta_2=0$; the alternative is that at least one coefficient is nonzero. There are $q=2$ restrictions.
\[
F=\frac{2.25+2.25-2(-0.50)(1.50)(1.50)}{2(1-0.25)}
=\frac{6.75}{1.50}=\boxed{4.50}.
\]
\answerpart{3}Reject the joint null, because $4.50>3.00$. The coefficient estimates move together in a direction that is relatively unusual given their negative estimated correlation.
\par\smallskip
The average $(t_1^2+t_2^2)/2=2.25$ ignores covariance. It equals the joint statistic only when $\hat\rho=0$.
\end{frame}
\fi

\begin{frame}{Problem 3: Are two study activities equally effective?}
\small
In a different large student sample, the estimated test-score slopes on weekly tutoring hours and independent-study hours are $\hat\beta_1=1.80$ and $\hat\beta_2=0.60$. Software supplies
\[
\widehat{\Var}(\hat\beta_1)=0.16,\text{ }\widehat{\Var}(\hat\beta_2)=0.09,
\text{ and }\widehat{\Cov}(\hat\beta_1,\hat\beta_2)=0.04.
\]
Use large-sample inference at 5\%.
\begin{enumerate}
\item Test $H_0:\beta_1=\beta_2$. Calculate the estimated difference, its standard error, and its $t$ statistic. How many restrictions does this null impose?
\item Construct a 95\% confidence interval for $\beta_1-\beta_2$. Calculate the equivalent $F$ statistic and compare it with the $q=1$ critical value $3.84$.
\item Rewrite $Y_i=\beta_0+\beta_1X_{1i}+\beta_2X_{2i}+U_i$ using $W_i=X_{1i}+X_{2i}$ so that one coefficient is $\beta_1-\beta_2$. What happens to the fitted values and $R^2$?
\end{enumerate}
\end{frame}

\ifshowsolutions
\begin{frame}{Solution 3: One restriction on two coefficients}
\small
\answerpart{1}Let $\delta=\beta_1-\beta_2$. Then $\hat\delta=1.20$ and
\[
SE(\hat\delta)=\sqrt{0.16+0.09-2(0.04)}=\sqrt{0.17}=0.4123.
\]
\[
t=\frac{1.20}{\sqrt{0.17}}=2.9104.
\]
There is one restriction, $\delta=0$. Reject equality at 5\%.
\par\smallskip
\answerpart{2}The 95\% interval is
\[
1.20\pm1.96\sqrt{0.17}=[0.3919,2.0081].
\]
It excludes zero. The equivalent statistic is $F=t^2=8.4706>3.84$, giving the same decision.
\par\smallskip
\answerpart{3}The equivalent regression is
\[
Y_i=\beta_0+(\beta_1-\beta_2)X_{1i}+\beta_2W_i+U_i.
\]
Test the coefficient on $X_1$ against zero. The regressors are rearranged without changing the fitted values, residuals, or $R^2$.
\end{frame}
\fi

\begin{frame}{Problem 4: Which coefficient pairs belong to the confidence set?}
\small
Reuse Problem 2: $(\hat\beta_1,\hat\beta_2)=(1.50,1.50)$, both standard errors are 1, and $\hat\rho=-0.50$.
For a candidate pair $(b_1,b_2)$, set $t_j=(\hat\beta_j-b_j)/SE(\hat\beta_j)$ and calculate
\[
F(b_1,b_2)=\frac{t_1^2+t_2^2-2\hat\rho t_1t_2}{2(1-\hat\rho^2)}.
\]
The approximate 95\% joint confidence set includes pairs with $F(b_1,b_2)\le3.00$.
\begin{enumerate}
\item Decide whether $(0,0)$, $(1,1)$, and $(0,1.50)$ belong to the joint set.
\item Construct the two separate 95\% intervals. Does their rectangular product contain $(0,0)$? Compare this with your joint decision.
\item In a separate illustration, suppose two individual 95\% coverage events are independent. Calculate the probability that both intervals cover, and the probability that at least one fails to cover. Why is the rectangle not automatically a 95\% joint set?
\end{enumerate}
\end{frame}

\ifshowsolutions
\begin{frame}{Solution 4: Inverting a joint test}
\small
\answerpart{1}Here $2(1-\hat\rho^2)=1.50$ and the numerator is $t_1^2+t_2^2+t_1t_2$.
\begin{center}
\begin{tabular}{cccc}
\toprule
Candidate pair & $(t_1,t_2)$ & $F$ & In the 95\% joint set?\\
\midrule
$(0,0)$ & $(1.50,1.50)$ & $4.50$ & No\\
$(1,1)$ & $(0.50,0.50)$ & $0.50$ & Yes\\
$(0,1.50)$ & $(1.50,0)$ & $1.50$ & Yes\\
\bottomrule
\end{tabular}
\end{center}
\answerpart{2}Each separate interval is $1.50\pm1.96=[-0.46,3.46]$.
Their rectangle contains $(0,0)$, although the joint test rejects that pair. Joint inference accounts for the covariance between estimates.
\par\smallskip
\answerpart{3}Under the independence assumption in the illustration,
\[
\Pr(\text{both cover})=0.95^2=0.9025,\text{ and }
\Pr(\text{at least one fails})=1-0.9025=0.0975.
\]
Separate 95\% intervals do not guarantee 95\% simultaneous coverage. These product calculations do not apply to the correlated estimates in Problem 2.
\end{frame}
\fi

\section{Specification}
\begin{frame}{Problem 5: Should we choose the model with the highest $R^2$?}
\small
Three regressions use the same 500 schools and test-score outcome. The coefficient of interest is on class size; supplied standard errors are in parentheses.
\begin{center}
\begin{tabular}{lccc}
\toprule
 & Model A & Model B & Model C\\
\midrule
Class-size coefficient & $-2.00$ & $-0.90$ & $-0.95$\\
Standard error & $(0.80)$ & $(0.40)$ & $(0.42)$\\
Income and parental education & No & Yes & Yes\\
Additional background controls & No & No & Yes\\
$R^2$ & $0.15$ & $0.45$ & $0.48$\\
\bottomrule
\end{tabular}
\end{center}
\begin{enumerate}
\item Lower-income districts tend to have larger classes and lower test scores for other reasons. Explain the possible omitted-variable bias in A and the movement from A to B.
\item Use a 5\% two-sided normal test for the class-size coefficient in each model. Does the largest $R^2$ establish the best causal estimate?
\item What does the comparison of B and C suggest about sensitivity to these extra controls? State what further reasoning is needed before calling the estimate causal.
\end{enumerate}
\end{frame}

\ifshowsolutions
\begin{frame}{Solution 5: Fit, significance, and causal interpretation}
\small
\answerpart{1}Income is positively related to scores and negatively related to class size in the stated setting. Omitting income can make the class-size coefficient too negative. Adding background controls makes the estimate less negative, from $-2.00$ to $-0.90$, consistent with this explanation.
\par\medskip
\answerpart{2}The three zero-null statistics are
\[
t_A=-2.00/0.80=-2.50,\text{ }t_B=-0.90/0.40=-2.25,
\text{ and }t_C=-0.95/0.42=-2.2619.
\]
All three reject at 5\%. Their significance does not establish causality, and a higher $R^2$ measures better in-sample fit rather than removal of confounding.
\par\medskip
\answerpart{3}The class-size estimate and its standard error change little from B to C. This suggests stability to these particular additional controls; it is not proof that all relevant confounders have been addressed.
\par\smallskip
Choose controls using the substantive question, their timing and role, and the plausibility of $\mathbb{E}[U_i\mid ClassSize_i,\text{controls}]=0$. Check sample selection and remaining omitted determinants before interpreting the association causally.
\end{frame}
\fi
\end{document}
